\subsection*{A3: Decision Making in Text World Environments} 
\label{sec:app:alfworld}
In this subsection, we demonstrate how \libName can be used to develop effective applications that involve interactive or online decision making. We perform the study using the ALFWorld~\cite{ALFWorld20} benchmark, which includes a diverse collection of synthetic language-based interactive decision-making tasks in household environments. 

With \libName, we implemented a two-agent system to solve tasks from  ALFWorld. It consists of an LLM-backed assistant agent responsible for suggesting plans to conduct a task and an executor agent responsible for executing actions in the ALFWorld environments. This system integrates ReAct prompting~\cite{yao2022react}, 
and is able to achieve similar performance. A common challenge encountered in both ReAct and the \libName-based two-agent system is their occasional inability to leverage basic commonsense knowledge about the physical world. This deficiency can lead to the system getting stuck in a loop due to repetitive errors. Fortunately, the modular design of \libName allows us to address this issue effectively: With \libName, we are able to introduce a grounding agent, which supplies crucial commonsense knowledge--such as \emph{“You must find and take the object before you can examine it. You must go to where the target object is before you can use it.”}--whenever the system exhibits early signs of recurring errors. It significantly enhances the system's ability to avoid getting entangled in error loops.
We compare the task-solving performance of the two variants of our system with GPT-3.5-turbo and ReAct\footnote{Results of ReAct are obtained by directly running its official code with default settings. The code uses \texttt{text-davinci-003} as backend LM and does not support GPT-3.5-turbo or GPT-4.} on the 134 unseen tasks from ALFWorld and report results in Figure~\ref{fig:res:alf}. The results show that introducing a grounding agent could bring in a 15\% performance gain on average. Upon examining the systems' outputs, we observe that the grounding agent, by delivering background commonsense knowledge at the right junctures, significantly mitigated the tendency of the system to persist with a flawed plan, thereby avoiding the creation of error loops. For an example trajectory comparing the systems see Appendix~\ref{appendix:app:alfworld}, Figure \ref{fig:alfchat_comparison}.

